
%%
%% This is file `sample-sigconf.tex',
%% generated with the docstrip utility.
%%
%% The original source files were:
%%
%% samples.dtx  (with options: `all,proceedings,bibtex,sigconf')
%% 
%% IMPORTANT NOTICE:
%% 
%% For the copyright see the source file.
%% 
%% Any modified versions of this file must be renamed
%% with new filenames distinct from sample-sigconf.tex.
%% 
%% For distribution of the original source see the terms
%% for copying and modification in the file samples.dtx.
%% 
%% This generated file may be distributed as long as the
%% original source files, as listed above, are part of the
%% same distribution. (The sources need not necessarily be
%% in the same archive or directory.)
%%
%%
%% Commands for TeXCount
%TC:macro \cite [option:text,text]
%TC:macro \citep [option:text,text]
%TC:macro \citet [option:text,text]
%TC:envir table 0 1
%TC:envir table* 0 1
%TC:envir tabular [ignore] word
%TC:envir displaymath 0 word
%TC:envir math 0 word
%TC:envir comment 0 0
%%
%% The first command in your LaTeX source must be the \documentclass
%% command.
%%
%% For submission and review of your manuscript please change the
%% command to \documentclass[manuscript, screen, review]{acmart}.
%%
%% When submitting camera ready or to TAPS, please change the command
%% to \documentclass[sigconf]{acmart} or whichever template is required
%% for your publication.
%%
%%
% VLDB template version of 2020-08-03 enhances the ACM template, version 1.7.0:
% https://www.acm.org/publications/proceedings-template
% The ACM Latex guide provides further information about the ACM template

\documentclass[sigconf, nonacm]{acmart}
\settopmatter{authorsperrow=5}


%% The following content must be adapted for the final version
% paper-specific
\newcommand\vldbdoi{XX.XX/XXX.XX}
\newcommand\vldbpages{XXX-XXX}
% issue-specific
\newcommand\vldbvolume{14}
\newcommand\vldbissue{1}
\newcommand\vldbyear{2020}
% should be fine as it is
\newcommand\vldbauthors{\authors}
\newcommand\vldbtitle{\shorttitle} 
% leave empty if no availability url should be set
\newcommand\vldbavailabilityurl{URL_TO_YOUR_ARTIFACTS}
% whether page numbers should be shown or not, use 'plain' for review versions, 'empty' for camera ready
\newcommand\vldbpagestyle{plain} 

%防止与已有包冲突
\makeatletter
\newif\if@restonecol
\makeatother
% \let\algorithm\relax
% \let\endalgorithm\relax
 
%引入伪代码模块需要的包
\usepackage[linesnumbered,ruled,vlined]{algorithm2e}
\SetKwComment{Comment}{// }{}
% \usepackage{algpseudocode}
\usepackage{amsmath}

\usepackage{enumitem}
\usepackage{multicol}
\usepackage{multirow}
\usepackage{float}
\usepackage{subfig}
\usepackage{balance}
\usepackage{bm}
\usepackage{tikz}
\usepackage{graphicx}

\newcommand{\xbr}[1]{\textcolor{teal}{\textbf{} #1}}

\newcommand{\circled}[2][0.8]{%
  % \raisebox{距离}{内容}：负值代表下沉，正值代表上升
  \raisebox{-0.2ex}{
    \begin{tikzpicture}
      \node[
        shape=circle, 
        draw, 
        inner sep=0pt,      % 清空内边距
        line width=0.4pt,   % 细边框保持精致
        minimum size=1.6ex % 稍微缩小尺寸，使其视觉高度与带升部字母（如 f）更协调
      ] (char) {\scalebox{#1}{#2}};
    \end{tikzpicture}%
  }%
}
\setlist[itemize]{leftmargin=*,nosep}
\setlist[enumerate]{leftmargin=*,nosep}
\setlength{\textfloatsep}{6pt plus 2pt minus 2pt}

\newcommand{\runin}[1]{\par\noindent\textbf{#1.} }
\newcommand{\subhead}[1]{\runin{#1}} % subhead非标准，弃用


    


\newcommand{\modora}{\texttt{MoDora}\xspace}
%%
%% \BibTeX command to typeset BibTeX logo in the docs
\AtBeginDocument{%
  \providecommand\BibTeX{{%
    Bib\TeX}}}

%% Rights management information.  This information is sent to you
%% when you complete the rights form.  These commands have SAMPLE
%% values in them; it is your responsibility as an author to replace
%% the commands and values with those provided to you when you
%% complete the rights form.
% \setcopyright{acmlicensed}
% \copyrightyear{2026}
% \acmYear{2026}
% \acmDOI{XXXXXXX.XXXXXXX}
%% These commands are for a PROCEEDINGS abstract or paper.
% \acmConference[SIGMOD/PODS '26]{International Conference on Management of Data}{May 31--June 05, 2026}{Bengaluru, India}
%%
%%  Uncomment \acmBooktitle if the title of the proceedings is different
%%  from ``Proceedings of ...''!
%%
%%\acmBooktitle{SIGMOD/PODS '26: International Conference on Management of Data,
%%  May 31--June 05, 2026, Bengaluru, India}
% \acmISBN{978-1-4503-XXXX-X/26/05}


%%
%% Submission ID.
%% Use this when submitting an article to a sponsored event. You'll
%% receive a unique submission ID from the organizers
%% of the event, and this ID should be used as the parameter to this command.
%%\acmSubmissionID{123-A56-BU3}

%%
%% For managing citations, it is recommended to use bibliography
%% files in BibTeX format.
%%
%% You can then either use BibTeX with the ACM-Reference-Format style,
%% or BibLaTeX with the acmnumeric or acmauthoryear sytles, that include
%% support for advanced citation of software artefact from the
%% biblatex-software package, also separately available on CTAN.
%%
%% Look at the sample-*-biblatex.tex files for templates showcasing
%% the biblatex styles.
%%

%%
%% The majority of ACM publications use numbered citations and
%% references.  The command \citestyle{authoryear} switches to the
%% "author year" style.
%%
%% If you are preparing content for an event
%% sponsored by ACM SIGGRAPH, you must use the "author year" style of
%% citations and references.
%% Uncommenting
%% the next command will enable that style.
%%\citestyle{acmauthoryear}


%%
%% end of the preamble, start of the body of the document source.
\begin{document}

%%
%% The "title" command has an optional parameter,
%% allowing the author to define a "short title" to be used in page headers.
\title{\modora: A Tree-Structured Multimodal RAG System for Complex Document Understanding}

%%
%% The "author" command and its associated commands are used to define the authors and their affiliations.
%%
%% The "author" command and its associated commands are used to define the authors and their affiliations.
\author{Yukai Wu}
\affiliation{%
  \institution{Shanghai Jiao Tong University}
  \city{}
  \country{}
}
\email{wuyukai99@gmail.com} % 换个gmail邮箱

\author{Bangrui Xu}
\affiliation{%
  \institution{Shanghai Jiao Tong University}
  \city{}
  \country{}
}
\email{dreamameter@sjtu.edu.cn}

\author{Zirui Tang}
\affiliation{
\institution{Shanghai Jiao Tong University}
\country{}
}
\email{tangzirui@sjtu.edu.cn}


\author{Xuanhe Zhou}
\authornote{Corresponding author}
\affiliation{%
  \institution{Shanghai Jiao Tong University}
  \city{}
  \country{}
}
\email{zhouxuanhe@sjtu.edu.cn}

\author{Fan Wu}
\affiliation{
    Shanghai Jiao Tong University
    \city{}
    \country{}
}
\email{fwu@cs.sjtu.edu.cn}


%%
%% By default, the full list of authors will be used in the page
%% headers. Often, this list is too long, and will overlap
%% other information printed in the page headers. This command allows
%% the author to define a more concise list
%% of authors' names for this purpose.
\renewcommand{\shortauthors}{Wu et al}

%%
%% The abstract is a short summary of the work to be presented in the
%% article.
\begin{abstract}
Traditional Retrieval-Augmented Generation (RAG) systems typically rely on a ``chunking + vector retrieval'' flat strategy, which frequently leads to the loss of hierarchical structure and visual context in complex multimodal documents. To address these limitations, we present \modora, a novel tree-structured multimodal RAG system designed for complex document understanding. \modora introduces a three-stage pipeline: (1) multimodal preprocessing to enrich non-textual elements using Vision Language Models (VLMs) into semantic triplets; (2) recursive tree construction to preserve logical hierarchies via a Component-Correlation Tree (CCTree); and (3) a two-pathway retrieval strategy that combines spatial location-based mapping with semantic ``Forward Search and Backward Verification.'' We demonstrate an interactive interface featuring a tree visualizer for manual structural refinement and interactive grounding for precise, verifiable evidence tracing within original PDF layouts. Experimental results on the MMDA benchmark show that \modora achieves an AIC-Acc of 71.1\%, outperforming the strongest baseline (M3DocRAG) by a margin of 23.7\%.
\end{abstract}

%%
%% The code below is generated by the tool at http://dl.acm.org/ccs.cfm.
%% Please copy and paste the code instead of the example below.
%%
% \begin{CCSXML}
% <ccs2012>
%  <concept>
%   <concept_id>00000000.0000000.0000000</concept_id>
%   <concept_desc>Do Not Use This Code, Generate the Correct Terms for Your Paper</concept_desc>
%   <concept_significance>500</concept_significance>
%  </concept>
%  <concept>
%   <concept_id>00000000.00000000.00000000</concept_id>
%   <concept_desc>Do Not Use This Code, Generate the Correct Terms for Your Paper</concept_desc>
%   <concept_significance>300</concept_significance>
%  </concept>
%  <concept>
%   <concept_id>00000000.00000000.00000000</concept_id>
%   <concept_desc>Do Not Use This Code, Generate the Correct Terms for Your Paper</concept_desc>
%   <concept_significance>100</concept_significance>
%  </concept>
%  <concept>
%   <concept_id>00000000.00000000.00000000</concept_id>
%   <concept_desc>Do Not Use This Code, Generate the Correct Terms for Your Paper</concept_desc>
%   <concept_significance>100</concept_significance>
%  </concept>
% </ccs2012>
% \end{CCSXML}

% \ccsdesc[500]{Do Not Use This Code~Generate the Correct Terms for Your Paper}
% \ccsdesc[300]{Do Not Use This Code~Generate the Correct Terms for Your Paper}
% \ccsdesc{Do Not Use This Code~Generate the Correct Terms for Your Paper}
% \ccsdesc[100]{Do Not Use This Code~Generate the Correct Terms for Your Paper}

%%
%% Keywords. The author(s) should pick words that accurately describe
%% the work being presented. Separate the keywords with commas.

% \keywords{Approximate nearest neighbor, Inverted lists, High-dimensional vectors}
  
%% A "teaser" image appears between the author and affiliation
%% information and the body of the document, and typically spans the
%% page.

%%
%% This command processes the author and affiliation and title
%% information and builds the first part of the formatted document.
\maketitle

%%% do not modify the following VLDB block %%
%%% VLDB block start %%%
% \pagestyle{\vldbpagestyle}
% \begingroup\small\noindent\raggedright\textbf{PVLDB Reference Format:}\\
% \vldbauthors. \vldbtitle. PVLDB, \vldbvolume(\vldbissue): \vldbpages, \vldbyear.\\
% \href{https://doi.org/\vldbdoi}{doi:\vldbdoi}
% \endgroup
% \begingroup
% \renewcommand\thefootnote{}\footnote{\noindent
% This work is licensed under the Creative Commons BY-NC-ND 4.0 International License. Visit \url{https://creativecommons.org/licenses/by-nc-nd/4.0/} to view a copy of this license. For any use beyond those covered by this license, obtain permission by emailing \href{mailto:info@vldb.org}{info@vldb.org}. Copyright is held by the owner/author(s). Publication rights licensed to the VLDB Endowment. \\
% \raggedright Proceedings of the VLDB Endowment, Vol. \vldbvolume, No. \vldbissue\ %
% ISSN 2150-8097. \\
% \href{https://doi.org/\vldbdoi}{doi:\vldbdoi} \\
% }\addtocounter{footnote}{-1}\endgroup
% %%% VLDB block end %%%

% %%% do not modify the following VLDB block %%
% %%% VLDB block start %%%
% \ifdefempty{\vldbavailabilityurl}{}{
% \vspace{.3cm}
% \begingroup\small\noindent\raggedright\textbf{PVLDB Artifact Availability:}\\
% The source code, data, and/or other artifacts have been made available at \url{\vldbavailabilityurl}.
% \endgroup
% }
%%% VLDB block end %%%

\section{Introduction}
\label{sec:intro}

% --- Introduction Figure ---
\begin{figure}[t]
    \centering
    \includegraphics[width=0.95\linewidth]{Fig/intro_v2.pdf}
    \caption{Overview of \modora with CCTree-based retrieval and interactive verification.}
    \label{fig:intro}
\end{figure}
% -----------------------------

Retrieval-Augmented Generation (RAG) has become a cornerstone for enhancing Large Language Models (LLMs) with external knowledge. By grounding responses in retrieved documents, RAG mitigates hallucinations and enables up-to-date reasoning. However, existing RAG pipelines typically rely on a ``chunking + vector retrieval'' flat strategy. This approach segments documents into independent text blocks, inevitably leading to the loss of hierarchical structure and visual context. While effective for plain text, this paradigm struggles with complex multimodal documents (e.g., financial reports, scientific papers) where layout and visual elements carry critical semantic information.

\runin{Challenges}
Despite recent advancements, three key challenges remain in multimodal document understanding. 
\textit{First, structural fragmentation.} Traditional RAG treats documents as a sequence of isolated chunks, failing to capture parent-child relationships between sections. This often leads to interpretations that are effectively ``taken out of context''. 
\textit{Second, semantic gaps in structured systems.} While systems like ZenDB~\cite{ZenDB} specialize in structured data, they often rely on pure text relations or SQL logic, ignoring non-textual visual cues. This can result in the retrieval of redundant or irrelevant evidence when queries involve charts or tables. 
\textit{Third, visual hallucinations in native VLMs.} End-to-end Vision Language Models (VLMs) lack explicit structural guidance. Without a logical hierarchy to navigate, they frequently misinterpret \xbr{document structure and connections between text and visual elements.}%data trends or lose the connection between legends and visual elements, especially in dense layouts.

\runin{Extension over Prior MoDora} The original MoDora work~\cite{xu2026modora} introduced CCTree-based semi-structured document understanding and released the MMDA benchmark. Building on that foundation, this paper emphasizes demo-oriented advances in the deployed backend: (1) an asynchronous OCR-to-CCTree pipeline with robust fallback for real-world documents, (2) tri-path evidence retrieval (location, semantic, vector) with unified grounding output, and (3) a verification-then-fallback QA loop with impact updates for interactive analysis.

\runin{Our Proposal} To address these limitations, we present \modora, a novel tree-structured multimodal RAG system designed for complex document understanding. As illustrated in Figure~\ref{fig:intro}, \modora constructs a hierarchical Component-\-Correlation Tree (CCTree) from unstructured PDFs. 
The system is built on three core designs: 
(1) \textit{Tree-Structured Modeling}, which preserves logical hierarchies and parent-child relationships instead of flat chunking; 
(2) \textit{VLM-Enhanced Retrieval}, which incorporates Vision Language Models to enrich non-textual elements into semantic triplets (Title, Metadata, Data), ensuring accurate alignment between visual evidence and textual queries; 
and (3) \textit{Interactive Verification}, which provides a frontend featuring \textit{Interactive Grounding} to trace answers back to original PDF regions, a \textit{Tree Visualizer} for manual or natural-language structural refinement, and a \textit{Node Heatmap} to identify critical information hubs. 
This human-in-the-loop design ensures that the retrieval process is not only accurate but also transparent and verifiable.

\runin{Results and Contributions} We have implemented \modora and evaluated it on the MMDA benchmark proposed in MoDora~\cite{xu2026modora}. Experimental results demonstrate that \modora achieves an AIC-Acc of 71.1\%, outperforming the strongest baseline (DocAgent~\cite{DocAgent}) by a margin of 14.1\%. 
In summary, our contributions are: 
(i) we propose a CCTree-based indexing structure that preserves document hierarchy and visual context for multimodal RAG; 
(ii) we design a query-adaptive retrieval strategy that coordinates location, semantic, and vector retrievers under a unified verification framework; 
and (iii) we demonstrate an interactive system for evidence grounding, node-impact analysis, and cross-session knowledge-base reuse.


\section{System Architecture}
\label{sec:architecture}

The architecture of \modora is designed to transform unstructured, multimodal documents into a structured knowledge hierarchy while providing an interactive and verifiable user experience. As illustrated in Figure~\ref{fig:arch}, the system is vertically organized into three layers: the \textbf{User Interaction Layer}, the \textbf{\modora Service Engine}, and the \textbf{Three-Stage Core Pipeline}.

\begin{figure*}[tp]
  \centering
  \includegraphics[width=0.9\linewidth]{Fig/arch(5).pdf}
  \caption{System Architecture of \modora.}
  \label{fig:arch}
\end{figure*}

\subsection{User Interaction Layer}
To bridge natural language queries with structured document trees, \modora provides an integrated environment featuring three specialized modules:

\runin{File Manager} Handles multimodal document ingestion and semantic tagging. \xbr{A persistent global library enables cross-session document reuse, eliminating redundant preprocessing.}%It maintains a persistent global library, enabling cross-session document reuse and eliminating redundant preprocessing.

\runin{Chat Engine} Supports complex QA with \textit{Interactive Grounding}. This feature allows users to trace AI-generated answers back to original PDF layouts via automated coordinate-based highlighting, ensuring the reasoning process is verifiable.

\runin{Tree Visualizer} Facilitates human-in-the-loop intervention by allowing users to manually refine the CCTree structure. It incorporates a \textit{Node Heatmap} (based on retrieval impact) to visually identify critical information hubs within the document hierarchy.

\subsection{\modora Service Engine}
Acting as the middleware bridging the user interaction layer and the core algorithmic pipeline, the Service Engine is optimized for multimodal RAG workloads. It orchestrates the asynchronous execution of computationally intensive VLM preprocessing tasks (Stage 1) and manages global document indexing to support cross-document question-answering across various sessions.

\subsection{Stage 1: Multimodal Preprocessing}
Before indexing, documents undergo a rigorous enrichment process to capture multimodal nuances:

\runin{Local-Alignment Aggregation} The system first performs optical character recognition (OCR) using \textit{PPStructureV3}~\cite{PPStructureV3} to identify raw elements such as titles, text blocks, tables, and images. Subsequently, a local-alignment strategy is applied using heuristic rules (e.g., text aggregation rule $A_1$, non-text aggregation rule $A_2$) to group fragmented elements into self-contained Aggregated Components. This ensures the components \xbr{are semantically (e.g., a table with its corresponding title) and visually (pixel-level coordinate) complete}.

\runin{VLM-based Information Enrichment} For non-textual components, \modora utilizes VLMs to extract integrated information, forming a semantic triplet consisting of Title, Metadata, and Data. For charts, the focus is on axes, labels, and trends; for tables, the schema and representative cell values are prioritized.

\subsection{Stage 2: Recursive Tree Construction}
The system organizes the enriched components into a logical hierarchy through a two-stage process:

\runin{Hierarchical Relationship Modeling} 1) \textit{Main Branch Construction}: Establishes ``text-to-text'' hierarchical relations based on title level prediction and ``text-to-other'' associations by attaching adjacent non-text components as children of relevant text nodes. 2) \textit{Supplementary Branch Isolation}: Page headers, footers, and sidebars are isolated into an independent branch under a virtual root to prevent semantic interference \xbr{with the main body} during retrieval.

\begin{sloppypar}
\runin{Bottom-up Metadata Summarization} To address information loss in long-document retrieval, summaries are recursively generated from leaf nodes upward. To balance information density, the number of summary keywords $n_i$ for each node $v_i$ is determined by an information attenuation formula, i.e., $n_i = \left\lceil \frac{D_i + d_i - 1}{D_i} \cdot \log_2(n_{i-1}^k) + 1 \right\rceil$, where $D_i$ denotes the maximum depth of the subtree rooted at $v_i$, $d_i$ represents the depth of node $v_i$, and $n_{i-1}$ indicates the total number of summaries generated by the children of $v_i$. The term $\frac{D_i + d_i - 1}{D_i}$ serves as the information attenuation coefficient, while $k$ is a hyperparameter controlling the rate of information growth. This ensures the root node captures a global overview while preserving critical details from the leaves.
\end{sloppypar}

\subsection{Stage 3: Two-Pathway Retrieval}
\modora uses a question-type-aware ``Parse-Retrieve-Verify'' workflow to balance recall and precision.

\runin{Query Parsing} The system first identifies whether a query contains spatial cues, then routes it to location-aware or semantic retrieval.

\runin{Location-Aware Path} For position-sensitive queries, \textit{LocationRetriever} traverses the CCTree with normalized bounding boxes and returns nodes overlapping the target regions, enabling direct grounding to layout coordinates.

\runin{Semantic Path} For content-centric queries, \textit{SemanticRetriever} recursively prunes irrelevant branches with LLM guidance, applies embedding fallback on pruned branches, and verifies candidate nodes using both text and cropped visual regions.

\runin{Verification-then-Fallback} Retrieved evidence is used for final reasoning. If consistency checks fail, the system falls back to whole-page reasoning. The retrieved nodes' impact scores are then updated for future tree optimization.


\section{Demonstration Scenarios}
\label{sec:demonstration}
\begin{figure*}[t] % 浣跨敤 figure* 鍙互璺ㄥ弻鏍忔樉绀猴紙濡傛灉鏄崟鏍忔ā鏉垮幓鎺?锛?  \centering
  \includegraphics[width=0.9\textwidth]{Fig/ui_v3.pdf} 
  \caption{A running example of \modora}
  \label{fig:gui}
\end{figure*}

This section demonstrates how \modora turns semi-structured documents into verifiable knowledge assets (Figure~\ref{fig:gui}).

\runin{Hierarchical Parsing and Inspection} After ingestion, \modora executes an asynchronous pipeline (\textit{OCR/fallback} $\rightarrow$ \textit{component aggregation} $\rightarrow$ \textit{tree construction} $\rightarrow$ \textit{metadata generation}) and stores intermediate artifacts (e.g., OCR blocks, component pack, tree JSON). Users inspect and refine this structure through the \textit{Tree Visualizer} (Figure~\ref{fig:gui}-\circled{2}).

\runin{Verified Multimodal QA} Users issue natural-language queries across documents, and \modora performs query-adaptive retrieval (location or semantic+vector). With \textit{Interactive Grounding} (Figure~\ref{fig:gui}-\circled{5}), citation clicks highlight source regions in the original PDF layout (Figure~\ref{fig:gui}-\circled{6}). Each evidence item keeps retriever provenance (location/semantic/vector), page index, and bounding boxes for transparent validation.

\runin{Human-in-the-Loop Optimization} The \textit{Node Heatmap} (Figure~\ref{fig:gui}-\circled{3}) visualizes retrieval impact, while users can refine the tree through direct node editing or natural-language restructuring instructions. Indexed CCTrees are archived in the \textit{Knowledge Base View} (Figure~\ref{fig:gui}-\circled{4}) for cross-session reuse, together with auto-generated document tags and structural statistics.


% \input{4-IVF-CLIP and HIVF-CLIP}

% \input{5-LSM-IVF}

\section{Experiments} \label{sec:experiments}

We conduct a series of evaluations on the MMDA benchmark to demonstrate the effectiveness of \modora in processing semi-structured documents.

\subsection{Experimental Configuration}
\runin{Dataset and Metric} Our evaluation uses the MMDA dataset introduced in MoDora~\cite{xu2026modora}, which contains 1,065 QA pairs in 537 documents with different layouts. We employ AIC-Acc (\%) as the primary evaluation metric, using an LLM-based judge to verify correctness of the responses while allowing for extraneous content.

\runin{Implementation Details} We implement a hybrid architecture. 
We deploy Qwen3-VL-8B-Instruct~\cite{Qwen3-VL} as the LLM for intermediate steps, including enrichment, hierarchy detection, question parsing and retrieval. We use the GPT-5 API~\cite{GPT-5} for the final high-level reasoning and answer generation.

\runin{Baselines} We compare \modora against nine state-of-the-art baselines: UDOP~\cite{UDOP}, DocOWL2~\cite{DocOwl2}, M3DocRAG~\cite{M3DocRAG}, SV-RAG~\cite{SV-RAG}, TextRAG~\cite{TextRAG}, QUEST~\cite{QUEST}, ZenDB~\cite{ZenDB}, DocAgent~\cite{DocAgent}, and GPT-5~\cite{GPT-5}.

\subsection{Performance Analysis} 
As illustrated in Figure~\ref{fig:exp1}, \modora achieves an AIC-Acc of 71.1\% on the MMDA benchmark.

\begin{figure}[h] 
\centering 
\includegraphics[width=0.95\linewidth]{Fig/exp_new.pdf} 
\caption{Experimental results on the MMDA dataset.} 
\label{fig:exp1} 
\end{figure}

\runin{Comparative Advantages} \modora outperforms the strongest baseline, DocAgent~\cite{DocAgent} (57.0\%), by a margin of 14.1\%. This improvement stems from the system's ability to transform fragmented OCR elements into layout-aware components, which are then organized into a hierarchical tree.

\runin{Structural and Hybrid Reasoning} While text-centric methods like TextRAG~\cite{TextRAG} (36.2\%) and ZenDB~\cite{ZenDB} (47.1\%) struggle with non-textual elements, \modora's tree-based retrieval accurately aligns semantic cues with visual data. Furthermore, our approach significantly exceeds the performance of end-to-end models such as DocOWL2~\cite{DocOwl2} (29.6\%) and UDOP~\cite{UDOP} (15.8\%), which often fail to capture fine-grained layout distinctions.

\runin{Noise Reduction} Compared to the vanilla GPT-5~\cite{GPT-5} (46.5\%), which processes full document pages, \modora's question-type-aware retrieval mechanism filters irrelevant content more effectively, resulting in higher reasoning precision.

\subsection{Case Study: Multi-Scenario Analysis}\label{subsec:case_study}
We present three representative cases from the MMDA benchmark to highlight how \modora handles complex queries that typically baffle traditional RAG systems.

\begin{table}[h]
\centering
\caption{Performance Analysis in Complex Scenarios}
\label{tab:case_study}
\resizebox{\columnwidth}{!}{%
\begin{tabular}{@{}lllc@{}}
\toprule
\textbf{Scenario} & \textbf{Challenge} & \textbf{Baseline Limitation} & \textbf{\modora} \\ \midrule
Hierarchy & Multi-level nesting & Context loss in flat chunking & \checkmark \\
\textit{(Sectional)} & Section dependencies & Misidentified parent-child links & \checkmark \\ \midrule
Hybrid & Text-Chart alignment & Hallucination in visual reasoning & \checkmark \\
\textit{(Cross-modal)} & Table-to-text links & Missing cross-element context & \checkmark \\ \midrule
Location & Spatial cues & Lack of coordinate metadata & \checkmark \\
\textit{(Visual-rich)} & Precise UI elements & Overlapping region filtering & \checkmark \\ \bottomrule
\end{tabular}%
}
\end{table}

\runin{Scenario 1: Structural Depth (C1)} In queries regarding nested subsections (e.g., Section III's second subsection), traditional flat-retrieval methods often fail because they lose the parent-child context. \modora utilizes the CCTree to accurately traverse the document hierarchy and retrieve complete logical branches.

\runin{Scenario 2: Multi-Modal Fusion (C2)} Hybrid questions requiring integrated reasoning over tables and text (e.g., matching a ``winter'' experiment ID to a specific table row) often lead to hallucinations in vanilla GPT-5 or text-only RAG. \modora's VLM-based enrichment generates semantic triplets that align visual cell data with textual descriptors.

\runin{Scenario 3: Spatial Precision (C3)} For queries like ``What is in the top-right logo?'', most baselines are unable to process locational metadata. \modora's Location-Based Path maps coordinates to grid-partitioned nodes, ensuring the retrieval of the exact visual element.


% \input{7-related work}

% \input{8-conclusion}

%%
%% The acknowledgments section is defined using the "acks" environment
%% (and NOT an unnumbered section). This ensures the proper
%% identification of the section in the article metadata, and the
%% consistent spelling of the heading.
% \begin{acks}
% To Robert, for the bagels and explaining CMYK and color spaces.
% \end{acks}

%%
%% The next two lines define the bibliography style to be used, and
%% the bibliography file.
\bibliographystyle{ACM-Reference-Format}
\balance
\bibliography{reference}


%%
%% If your work has an appendix, this is the place to put it.
% \appendix

% \section{Research Methods}

% \subsection{Part One}

% Lorem ipsum dolor sit amet, consectetur adipiscing elit. Morbi
% malesuada, quam in pulvinar varius, metus nunc fermentum urna, id
% sollicitudin purus odio sit amet enim. Aliquam ullamcorper eu ipsum
% vel mollis. Curabitur quis dictum nisl. Phasellus vel semper risus, et
% lacinia dolor. Integer ultricies commodo sem nec semper.

% \subsection{Part Two}

% Etiam commodo feugiat nisl pulvinar pellentesque. Etiam auctor sodales
% ligula, non varius nibh pulvinar semper. Suspendisse nec lectus non
% ipsum convallis congue hendrerit vitae sapien. Donec at laoreet
% eros. Vivamus non purus placerat, scelerisque diam eu, cursus
% ante. Etiam aliquam tortor auctor efficitur mattis.

% \section{Online Resources}

% Nam id fermentum dui. Suspendisse sagittis tortor a nulla mollis, in
% pulvinar ex pretium. Sed interdum orci quis metus euismod, et sagittis
% enim maximus. Vestibulum gravida massa ut felis suscipit
% congue. Quisque mattis elit a risus ultrices commodo venenatis eget
% dui. Etiam sagittis eleifend elementum.

% Nam interdum magna at lectus dignissim, ac dignissim lorem
% rhoncus. Maecenas eu arcu ac neque placerat aliquam. Nunc pulvinar
% massa et mattis lacinia.

\end{document}
\endinput
%%
%% End of file `sample-sigconf.tex'.
